% Kalman (1960) 中文校订版父级 LaTeX 导言区配置
\usepackage[margin=2.4cm]{geometry}
\setlength{\emergencystretch}{2em}
\usepackage{amsmath,amssymb,amsthm,mathrsfs,bm}
\usepackage{booktabs,tabularx,array}
\usepackage{pifont,newunicodechar}
\usepackage{xcolor}
\definecolor{accent}{HTML}{1F6FB2}
\usepackage{tikz}
\usepackage{standalone}
\usepackage{etoolbox}

% Kalman (1960) 图 1–4 共享样式（原生适配 html-tex-doc 深浅色自适应规范色）。
% 用法：\documentclass[tikz,border=4pt]{standalone} 后 \input{common} 再 \input{kalman-style}。
\usepackage{bm,pifont}
\usetikzlibrary{positioning,fit,calc,arrows.meta,backgrounds,shapes.geometric}

% 最优解空心亮星上标（与正文 ✫ 一致）
\newcommand{\optstar}{\text{\ding{75}}}

% 深浅色自适应规范色（配合 htd/svgtheme.py）：
% black -> currentColor；black!6 -> currentColor (opacity≈0.059)；
% white -> var(--bg)；accent (#1F6FB2) -> var(--accent)。
\colorlet{edge}{black}
\colorlet{blk}{black!6}

\tikzset{
  kalman/.style={font=\small, >={Stealth[length=2.4mm]}, line width=0.7pt, draw=edge},
  block/.style={draw=edge, fill=blk, rounded corners=1.5pt,
                minimum height=8mm, minimum width=13mm, inner sep=3pt, text=edge},
  sum/.style={draw=edge, fill=white, circle, minimum size=6.5mm, inner sep=0pt, text=edge},
  tri/.style={draw=edge, fill=blk, isosceles triangle, isosceles triangle apex angle=60,
              shape border rotate=0, inner sep=1pt, minimum height=13mm,
              font=\footnotesize, align=center, text=edge},
  sig/.style={->, edge},
  lbl/.style={font=\footnotesize, text=edge},
  dot/.style={circle, fill=edge, inner sep=1.1pt},
  frame/.style={draw=accent, dashed, rounded corners=3pt, line width=0.6pt},
  framelbl/.style={lbl, text=accent, font=\scriptsize, anchor=south},
}

% 求和节点两个输入侧 (左/下) 的 + 号，置于圆外
\newcommand{\sumsigns}[1]{%
  \node[font=\scriptsize, text=edge] at ($(#1.west)+(-0.12,0.16)$) {$+$};
  \node[font=\scriptsize, text=edge] at ($(#1.south)+(-0.16,-0.12)$) {$+$};}

\newcommand{\TXTdelay}{单位\\延迟}
\newcommand{\TXTmodel}{随机过程模型}
\newcommand{\TXTplant}{被控物理系统}

\newunicodechar{✫}{\ding{75}}

\usepackage[colorlinks=true,linkcolor=accent,citecolor=accent,urlcolor=accent,
            pdftitle={线性滤波与预测问题的一种新方法},
            pdfauthor={R. E. Kalman}]{hyperref}

\renewcommand{\proofname}{证明}

\theoremstyle{plain}
\newtheorem{theorem}{定理}
\newtheorem{corollary}{推论}
\newtheorem{htdtheoremn}{定理}

\theoremstyle{definition}
\newtheorem{example}{算例}
\newtheorem{htdproblemn}{问题}
\newtheorem{htdremarkn}{评注}

\NewDocumentEnvironment{theoremn}{m o}{%
  \renewcommand{\thehtdtheoremn}{#1}%
  \IfValueTF{#2}{\begin{htdtheoremn}[#2]}{\begin{htdtheoremn}}%
}{\end{htdtheoremn}}

\NewDocumentEnvironment{problemn}{m o}{%
  \renewcommand{\thehtdproblemn}{#1}%
  \IfValueTF{#2}{\begin{htdproblemn}[#2]}{\begin{htdproblemn}}%
}{\end{htdproblemn}}

\NewDocumentEnvironment{remarkn}{m o}{%
  \renewcommand{\thehtdremarkn}{#1}%
  \IfValueTF{#2}{\begin{htdremarkn}[#2]}{\begin{htdremarkn}}%
}{\end{htdremarkn}}

\makeatletter
\newsavebox{\htd@figbox}
\newcommand{\htdfig}[3]{%
  \begin{figure}[htbp]
    \centering
    \filename@parse{html/src/#1}%
    \begin{lrbox}{\htd@figbox}%
      \includestandalone{\filename@area\filename@base}%
    \end{lrbox}%
    \ifdim\wd\htd@figbox>\linewidth
      \resizebox{\linewidth}{!}{\usebox{\htd@figbox}}%
    \else
      \usebox{\htd@figbox}%
    \fi
    \caption{#3}\label{#2}%
  \end{figure}%
}

\patchcmd{\thebibliography}{\section*{\refname}}{}{}{}
\makeatother

\newenvironment{htdtable}[1]{%
  \par\medskip\noindent
  \renewcommand{\arraystretch}{1.25}%
  \tabularx{\textwidth}{#1}%
}{%
  \endtabularx
  \par\medskip
}

\DeclareRobustCommand{\htdfnmark}[1]{%
  \hypertarget{htdfnref:#1}{}\textsuperscript{\hyperlink{htdfn:#1}{[#1]}}%
}
\newcommand{\htdfntext}[2]{%
  \par\smallskip\noindent
  \hypertarget{htdfn:#1}{\hyperlink{htdfnref:#1}{\textsuperscript{[#1]}}}~{\small #2}\par
}

\bibliographystyle{unsrt}
\AtBeginDocument{%
  \nocite{r01,r02,r03,r04,r05,r06,r07,r08,r09,r10,r11,r12,r13,r14,r15,%
          r16,r17,r18,r19,r20,r21,r22,r23,r24,r25,r26,r27,r28,r29}%
}
